% ============================================================
\setpartpreamble{\partintro
  {How a plugin is configured, registers its services and connects to the application's DI
  container.}
  {
    \item Chapter~\ref{ch:the-options-extension}: The options extension
    \item Chapter~\ref{ch:registering-services}: Registering services
    \item Chapter~\ref{ch:public-api-and-integration-into}: Public API and integration into
    the application
  }}
\part{Infrastructure}\label{part:II}

% ------------------------------------------------------------
\chapter{The options extension}\label{ch:the-options-extension}

The entry point of every real plugin is a class implementing
\code{IDb\allowbreak{}Context\allowbreak{}Options\allowbreak{}Extension}. It has three jobs:
it carries the configuration, it registers the services in the internal provider in
\code{Apply\allowbreak{}Services}, and it describes itself to the cache through \code{Info}.

\section{Configuration as a record}

All model-relevant settings of the sample live in a record. Records provide \code{Equals}
and \code{Get\allowbreak{}Hash\allowbreak{}Code} over all fields, which is exactly what the
cache key needs. New settings are added to the record without touching hashing or
comparison.

\begin{sourcecode}[Infrastructure/\allowbreak{}Tenancy\allowbreak{}Settings.\allowbreak{}cs]{csharp}
public sealed record TenancySettings
{
    public static TenancySettings Default { get; } = new();

    public string SettingName { get; init; } = "app.tenant_id";
    public string PolicyName { get; init; } = "tenant_isolation";
    public bool ForceRowLevelSecurity { get; init; } = true;
    public bool ApplyQueryFilter { get; init; }
    public bool CreateTenantIndex { get; init; } = true;
}
\end{sourcecode}

\section{Immutability}

EF compares and caches options. Once added, an extension must therefore never change again;
every change creates a copy. The pattern from EF Core itself: a protected copy constructor
and a virtual \code{Clone} method, so derived extensions clone their own type.

\begin{sourcecode}[Infrastructure/\allowbreak{}Tenancy\allowbreak{}Options\allowbreak{}Extension.\allowbreak{}cs (\allowbreak{}excerpt)]{csharp}
public class TenancyOptionsExtension : IDbContextOptionsExtension
{
    private DbContextOptionsExtensionInfo? _info;

    public TenancyOptionsExtension()
    {
    }

    protected TenancyOptionsExtension(TenancyOptionsExtension copyFrom)
    {
        ArgumentNullException.ThrowIfNull(copyFrom);

        Settings = copyFrom.Settings;
    }

    public virtual TenancySettings Settings { get; private set; } = TenancySettings.Default;

    public virtual DbContextOptionsExtensionInfo Info => _info ??= new ExtensionInfo(this);

    protected virtual TenancyOptionsExtension Clone() => new(this);

    public virtual TenancyOptionsExtension WithSettings(TenancySettings settings)
    {
        var clone = Clone();
        clone.Settings = settings;
        return clone;
    }
\end{sourcecode}

\section{The info class and the service provider cache}

The info class decides whether EF may reuse a cached internal service provider. Two methods
must agree for this:
\code{Get\allowbreak{}Service\allowbreak{}Provider\allowbreak{}Hash\allowbreak{}Code} as a
fast pre-filter and
\code{Should\allowbreak{}Use\allowbreak{}Same\allowbreak{}Service\allowbreak{}Provider} as
the exact comparison.

\begin{sourcecode}[Infrastructure/\allowbreak{}Tenancy\allowbreak{}Options\allowbreak{}Extension.\allowbreak{}cs › Extension\allowbreak{}Info]{csharp}
private sealed class ExtensionInfo : DbContextOptionsExtensionInfo
{
    public ExtensionInfo(TenancyOptionsExtension extension)
        : base(extension)
    {
    }

    private new TenancyOptionsExtension Extension => (TenancyOptionsExtension)base.Extension;

    public override bool IsDatabaseProvider => false;

    public override string LogFragment
        => $"AcmeTenancy(Setting={Extension.Settings.SettingName}) ";

    public override int GetServiceProviderHashCode() => Extension.Settings.GetHashCode();

    public override bool ShouldUseSameServiceProvider(DbContextOptionsExtensionInfo other)
        => other is ExtensionInfo o && o.Extension.Settings == Extension.Settings;

    public override void PopulateDebugInfo(IDictionary<string, string> debugInfo)
    {
        debugInfo["Acme:Tenancy:Force"] =
            Extension.Settings.ForceRowLevelSecurity.ToString();
        // ... further settings
    }
}
\end{sourcecode}

There are two classic mistakes with these methods, with opposite consequences:

\begin{datatable}{H{0.65}H{1.35}}
  \textbf{Mistake} & \textbf{Consequence} \\
  \tablehead
  A model-relevant option is missing from the hash & Two contexts with different
  configuration share provider and model. The second gets the cached, \textbf{wrong} model.
  Often unnoticed in tests, because only one configuration runs there. \\
  \tablerule
  A value that varies per context is in the hash (tenant, connection string, instance) & A
  new internal provider including model is created for every context. Memory and startup
  time explode; EF warns after 20 providers with
  \code{Many\allowbreak{}Service\allowbreak{}Providers\allowbreak{}Created\allowbreak{}Warning}.
  \\
\end{datatable}

\begin{important}[The rule of thumb for the hash]
  The hash must contain exactly what \textbf{influences which services are registered or
  what the model looks like}, and nothing else. The tenant itself changes per request and
  must therefore never end up in the options; it arrives at run time through an interceptor
  (chapter~\ref{ch:interceptors}).
\end{important}

\code{Log\allowbreak{}Fragment} appears in EF's startup message (“initialized using provider
… with options: …”) and helps when troubleshooting.
\code{Populate\allowbreak{}Debug\allowbreak{}Info} supplies values for EF's diagnostics when
too many service providers are created. \code{Is\allowbreak{}Database\allowbreak{}Provider}
is \code{true} only for real providers.

\section{Validation}

\code{Validate} is called when the service provider is created. Here the plugin checks
preconditions arising from the combination of extensions and reports them with a clear
message instead of a later follow-up error:

\begin{sourcecode}[Infrastructure/\allowbreak{}Tenancy\allowbreak{}Options\allowbreak{}Extension.\allowbreak{}cs › Validate]{csharp}
public virtual void Validate(IDbContextOptions options)
{
    if (options.FindExtension<NpgsqlOptionsExtension>() is null)
        throw new InvalidOperationException(
            "UseAcmeTenancy() requires UseNpgsql(). Other databases need their own "
            + "provider package.");
}
\end{sourcecode}

% ------------------------------------------------------------
\chapter{Registering services}\label{ch:registering-services}

\code{Apply\allowbreak{}Services} fills the internal service provider. EF provides a builder
for this that knows the lifetime and cardinality of every EF service:
\code{Entity\allowbreak{}Framework\allowbreak{}Services\allowbreak{}Builder} for core
services and
\code{Entity\allowbreak{}Framework\allowbreak{}Relational\allowbreak{}Services\allowbreak{}Builder}
additionally for relational services. It registers each service with the correct lifetime
and automatically uses \code{Try\allowbreak{}Add\allowbreak{}Enumerable} for services that
can be registered multiple times.

\begin{sourcecode}[Infrastructure/\allowbreak{}Tenancy\allowbreak{}Options\allowbreak{}Extension.\allowbreak{}cs › Apply\allowbreak{}Services]{csharp}
public virtual void ApplyServices(IServiceCollection services)
{
    // 1) Plugin interfaces: multiple registrations, composable with other plugins
    new EntityFrameworkRelationalServicesBuilder(services)
        .TryAdd<IRelationalTypeMappingSourcePlugin, TenancyTypeMappingSourcePlugin>()
        .TryAdd<IMethodCallTranslatorPlugin, TenancyMethodCallTranslatorPlugin>()
        .TryAdd<IMemberTranslatorPlugin, TenancyMemberTranslatorPlugin>()
        .TryAdd<IEvaluatableExpressionFilterPlugin,
                TenancyEvaluatableExpressionFilterPlugin>()
        .TryAdd<IConventionSetPlugin, TenancyConventionSetPlugin>()
        .TryAdd<ISingletonOptions, ITenancySingletonOptions>(
            p => p.GetRequiredService<ITenancySingletonOptions>())
        // 2) Own services that EF does not know
        .TryAddProviderSpecificServices(m => m
            .TryAddSingleton<ITenancySingletonOptions, TenancySingletonOptions>()
            .TryAddSingleton<TenancyLoggingDefinitions, TenancyLoggingDefinitions>());

    // 3) Replace single services of the provider
    services.Replace(ServiceDescriptor
        .Scoped<IMigrationsSqlGenerator, TenancyMigrationsSqlGenerator>());
    services.Replace(ServiceDescriptor
        .Singleton<IRelationalAnnotationProvider, TenancyAnnotationProvider>());
}
\end{sourcecode}

\begin{datatable}{H{0.83}H{0.77}H{1.40}}
  \textbf{Method} & \textbf{When} & \textbf{Why} \\
  \tablehead
  \code{Try\allowbreak{}Add<\allowbreak{}TService,\allowbreak{} TImpl>(\allowbreak{})} & EF
  services and plugin interfaces & Lifetime and multiple registration come from EF's table.
  Wrong lifetimes otherwise cause exceptions or subtly wrong behavior. \\
  \tablerule
  \code{Try\allowbreak{}Add\allowbreak{}Provider\allowbreak{}Specific\allowbreak{}Services}
  & Own services EF does not know & The builder rejects unknown types; here you set the
  lifetime yourself. \\
  \tablerule
  \code{services.\allowbreak{}Replace(\allowbreak{}.\allowbreak{}.\allowbreak{}.\allowbreak{})}
  & Replace a single provider service & The provider registers with
  \code{Try\allowbreak{}Add}. \code{Replace} wins regardless of whether the plugin runs
  before or after the provider. \\
\end{datatable}

\begin{note}[Why Replace works in both orders]
  EF calls the extensions' \code{Apply\allowbreak{}Services} in the order they were added.
  If the plugin runs first, \code{Replace} adds the service and the provider's later
  \code{Try\allowbreak{}Add} does nothing. If it runs afterwards, \code{Replace} replaces
  the provider's service. Both orders lead to the same result.
\end{note}

\section{Singleton options}

Singleton services such as type mapping plugins cannot resolve
\code{IDb\allowbreak{}Context\allowbreak{}Options} per context. For them there is
\code{ISingleton\allowbreak{}Options}: a singleton that EF initializes once from the options
when creating the provider and validates for every further context.

\begin{sourcecode}[Infrastructure/\allowbreak{}ITenancy\allowbreak{}Singleton\allowbreak{}Options.\allowbreak{}cs]{csharp}
public class TenancySingletonOptions : ITenancySingletonOptions
{
    public TenancySettings Settings { get; private set; } = TenancySettings.Default;

    // Once per internal service provider
    public virtual void Initialize(IDbContextOptions options)
        => Settings = options.FindExtension<TenancyOptionsExtension>()?.Settings
                      ?? TenancySettings.Default;

    // For every context: options must not differ from the provider's state
    public virtual void Validate(IDbContextOptions options)
    {
        var settings = options.FindExtension<TenancyOptionsExtension>()?.Settings
                       ?? TenancySettings.Default;
        if (settings != Settings)
            throw new InvalidOperationException("Tenancy options changed afterwards ...");
    }
}
\end{sourcecode}

The double registration in \code{Apply\allowbreak{}Services} is required: once under its own
interface so services can inject it, and once as \code{ISingleton\allowbreak{}Options} so EF
finds and initializes it. The factory variant ensures that both registrations return
\emph{the same} instance.

\section{Staying composable despite replacement}

When a plugin replaces a single service, a second plugin doing the same cannot be active at
the same time. Professional plugins mitigate this by additionally offering the actual logic
publicly and reusably. In the sample, \code{Create\allowbreak{}Policy\allowbreak{}Sql} is
public and all generator methods are virtual. A user who must combine two plugins writes
their own generator deriving from one and calling the logic of the other. This option
belongs in the documentation.

Better still, a plugin avoids the conflict altogether by \textbf{decorating} the service
instead of replacing it: it takes whatever implementation is registered at that point (the
provider's or another plugin's), wraps it and delegates to it. An example is a (fictional)
audit plugin that attaches a trigger to every new table. In its
\code{Apply\allowbreak{}Services} it looks up the existing registration and replaces only
the descriptor, not the logic:

\begin{sourcecode}[tests/\allowbreak{}Acme.\allowbreak{}Entity\allowbreak{}Framework\allowbreak{}Core.\allowbreak{}Tenancy.\allowbreak{}Tests/\allowbreak{}Composition\allowbreak{}Tests.\allowbreak{}cs › Audit\allowbreak{}Options\allowbreak{}Extension]{csharp}
public void ApplyServices(IServiceCollection services)
{
    // Registered so far: the provider's generator or that of another plugin
    var inner = services.LastOrDefault(d => d.ServiceType == typeof(IMigrationsSqlGenerator))
        ?? throw new InvalidOperationException("Call UseContosoAudit() after UseNpgsql().");

    services.Replace(ServiceDescriptor.Describe(
        typeof(IMigrationsSqlGenerator),
        sp => new AuditMigrationsSqlGenerator(
            (IMigrationsSqlGenerator)CreateInstance(sp, inner),
            sp.GetRequiredService<MigrationsSqlGeneratorDependencies>()),
        inner.Lifetime));     // lifetime of the original
}

// Creates the implementation the way the original descriptor would have
private static object CreateInstance(IServiceProvider sp, ServiceDescriptor descriptor)
    => descriptor.ImplementationInstance
       ?? descriptor.ImplementationFactory?.Invoke(sp)
       ?? ActivatorUtilities.CreateInstance(sp, descriptor.ImplementationType!);
\end{sourcecode}

The trick works because EF calls \code{Apply\allowbreak{}Services} in the order in which the
extensions were added; when the audit plugin runs, the provider has already registered its
generator. \code{Create\allowbreak{}Instance} reproduces the original registration exactly:
an instance, a factory, or a type whose dependencies (such as Npgsql's singleton options)
\code{Activator\allowbreak{}Utilities} resolves from the container. The decorator itself
only uses the public interface:

\begin{sourcecode}[tests/\allowbreak{}.\allowbreak{}.\allowbreak{}.\allowbreak{}/\allowbreak{}Composition\allowbreak{}Tests.\allowbreak{}cs › Audit\allowbreak{}Migrations\allowbreak{}Sql\allowbreak{}Generator (\allowbreak{}shortened)]{csharp}
public class AuditMigrationsSqlGenerator : IMigrationsSqlGenerator
{
    // Constructor: (IMigrationsSqlGenerator inner, MigrationsSqlGeneratorDependencies deps)

    public IReadOnlyList<MigrationCommand> Generate(
        IReadOnlyList<MigrationOperation> operations,
        IModel? model = null,
        MigrationsSqlGenerationOptions options = MigrationsSqlGenerationOptions.Default)
    {
        var commands = _inner.Generate(operations, model, options).ToList();   // unchanged
        var builder = new MigrationCommandListBuilder(_dependencies);
        var sql = _dependencies.SqlGenerationHelper;

        // Appended at the end: by then every table of this migration exists
        foreach (var table in operations.OfType<CreateTableOperation>())
        {
            builder.Append("CREATE TRIGGER audit_changes ")
                .Append("AFTER INSERT OR UPDATE OR DELETE ON ")
                .Append(sql.DelimitIdentifier(table.Name, table.Schema))
                .Append(" FOR EACH ROW EXECUTE FUNCTION audit.log_change()")
                .AppendLine(sql.StatementTerminator)
                .EndCommand();
        }

        commands.AddRange(builder.GetCommandList());
        return commands;
    }
}
\end{sourcecode}

In the application, the two plugins are combined without a single line of glue code:

\begin{sourcecode}[Program.\allowbreak{}cs]{csharp}
services.AddDbContext<ShopContext>((sp, options) => options
    .UseNpgsql(connectionString)
    .UseAcmeTenancy(sp)       // replaces IMigrationsSqlGenerator
    .UseContosoAudit());      // wraps whatever is registered: here Acme's
\end{sourcecode}

\begin{important}[Decorate where possible, replace where necessary]
  Decorators compose with each other in any order, because each one wraps whatever it finds.
  A plugin that uses \code{Replace} however discards everything registered before it,
  including decorators. A replacing plugin therefore has to be activated first, and a
  decorator must check that there is something to decorate and fail with a clear message
  otherwise (the provider is missing or was added later).

  Decorating requires that the public interface is enough. The tenancy plugin needs
  protected hooks inside \code{CREATE TABLE} and therefore derives from Npgsql's generator.
  For such cases the fallback described above remains: public logic such as
  \code{Create\allowbreak{}Policy\allowbreak{}Sql}, virtual methods, and
  \code{Replace\allowbreak{}Service} in the application.
\end{important}

\begin{tip}
  The test \code{Composition\allowbreak{}Tests} in the sample contains the complete audit
  plugin. It checks that the decorator wraps Acme's generator, that the migration contains
  the SQL of both plugins, and that activating the decorator before the provider fails.
\end{tip}

% ------------------------------------------------------------
\chapter{Public API and integration into the application}\label{ch:public-api-and-integration-into}

\section{The Use method}

Users activate the plugin through an extension method on
\code{Db\allowbreak{}Context\allowbreak{}Options\allowbreak{}Builder}. It fetches the
existing extension or creates a new one and registers it through
\code{IDb\allowbreak{}Context\allowbreak{}Options\allowbreak{}Builder\allowbreak{}Infrastructure}:

\begin{sourcecode}[Infrastructure/\allowbreak{}Tenancy\allowbreak{}Db\allowbreak{}Context\allowbreak{}Options\allowbreak{}Builder\allowbreak{}Extensions.\allowbreak{}cs (\allowbreak{}excerpt)]{csharp}
public static DbContextOptionsBuilder UseAcmeTenancy(
    this DbContextOptionsBuilder optionsBuilder,
    Action<TenancyOptionsBuilder>? configure = null)
{
    ArgumentNullException.ThrowIfNull(optionsBuilder);

    var extension = optionsBuilder.Options.FindExtension<TenancyOptionsExtension>()
                    ?? new TenancyOptionsExtension();

    ((IDbContextOptionsBuilderInfrastructure)optionsBuilder).AddOrUpdateExtension(extension);
    configure?.Invoke(new TenancyOptionsBuilder(optionsBuilder));
    return optionsBuilder;
}

// Generic overload: method chaining stays typed
public static DbContextOptionsBuilder<TContext> UseAcmeTenancy<TContext>(
    this DbContextOptionsBuilder<TContext> optionsBuilder,
    Action<TenancyOptionsBuilder>? configure = null) where TContext : DbContext
    => (DbContextOptionsBuilder<TContext>)UseAcmeTenancy(
        (DbContextOptionsBuilder)optionsBuilder, configure);
\end{sourcecode}

The file lives in the namespace
\code{Microsoft.\allowbreak{}Entity\allowbreak{}Framework\allowbreak{}Core}. That way the
method can be found without an additional \code{using}, just like
\code{Use\allowbreak{}Npgsql}. This is the common convention in the EF ecosystem.

\section{The options builder}

Settings are configured through a dedicated builder. It validates input immediately and
writes every change back as a new extension:

\begin{sourcecode}[Infrastructure/\allowbreak{}Tenancy\allowbreak{}Options\allowbreak{}Builder.\allowbreak{}cs (\allowbreak{}excerpt)]{csharp}
public class TenancyOptionsBuilder
{
    public TenancyOptionsBuilder(DbContextOptionsBuilder optionsBuilder)
    {
        ArgumentNullException.ThrowIfNull(optionsBuilder);

        OptionsBuilder = optionsBuilder;
    }

    protected virtual DbContextOptionsBuilder OptionsBuilder { get; }

    public virtual TenancyOptionsBuilder UseSettingName(string settingName)
        => With(s => s with
        {
            SettingName = TenancyIdentifiers.ValidateSettingName(settingName),
        });

    public virtual TenancyOptionsBuilder ApplyQueryFilter(bool apply = true)
        => With(s => s with { ApplyQueryFilter = apply });

    protected virtual TenancyOptionsBuilder With(
        Func<TenancySettings, TenancySettings> change)
    {
        var extension = OptionsBuilder.Options.FindExtension<TenancyOptionsExtension>()
                        ?? new TenancyOptionsExtension();

        ((IDbContextOptionsBuilderInfrastructure)OptionsBuilder)
            .AddOrUpdateExtension(extension.WithSettings(change(extension.Settings)));
        return this;
    }
}
\end{sourcecode}

\begin{note}[Alternative: activation through the provider builder]
  Provider-specific plugins are sometimes activated through the provider's options builder
  instead, for example \code{Use\allowbreak{}Npgsql(\allowbreak{}cs,\allowbreak{} o =>
  o.\allowbreak{}Use\allowbreak{}Noda\allowbreak{}Time(\allowbreak{}))}. For this you write
  an extension method on
  \code{Npgsql\allowbreak{}Db\allowbreak{}Context\allowbreak{}Options\allowbreak{}Builder}
  and reach the actual builder through
  \code{IRelational\allowbreak{}Db\allowbreak{}Context\allowbreak{}Options\allowbreak{}Builder\allowbreak{}Infrastructure.\allowbreak{}Options\allowbreak{}Builder}.
  The PostgreSQL package uses the simpler route via
  \code{Db\allowbreak{}Context\allowbreak{}Options\allowbreak{}Builder}; the SQL Server
  package in chapter~\ref{ch:a-second-provider-sql-server} shows the provider-builder style.
\end{note}

\section{Connecting two containers}

Some parts of a plugin need services of the \emph{application}: the connection interceptor
needs the current tenant, which only the application knows. The internal EF provider cannot
resolve app services, though. The clean solution: such components live in the app container
and are handed to EF via \code{Add\allowbreak{}Interceptors}.

\begin{sourcecode}[Dependency\allowbreak{}Injection/\allowbreak{}… (\allowbreak{}core),\allowbreak{} Infrastructure/\allowbreak{}… (\allowbreak{}Postgre\allowbreak{}SQL),\allowbreak{} excerpts]{csharp}
// Core package: only register the tenant provider in the app container
public static IServiceCollection AddAcmeTenancy<TTenantProvider>(
    this IServiceCollection services,
    ServiceLifetime lifetime = ServiceLifetime.Scoped)
    where TTenantProvider : class, ITenantProvider
{
    services.TryAdd(ServiceDescriptor.Describe(
        typeof(ITenantProvider), typeof(TTenantProvider), lifetime));
    return services;
}

// Provider package: Use overload creates the interceptors from this provider
public static DbContextOptionsBuilder UseAcmeTenancy(
    this DbContextOptionsBuilder optionsBuilder,
    IServiceProvider applicationServices,
    Action<TenancyOptionsBuilder>? configure = null)
{
    var tenants = applicationServices.GetRequiredService<ITenantProvider>();

    return optionsBuilder
        .UseAcmeTenancy(configure)
        .AddInterceptors(
            new TenantConnectionInterceptor(tenants),
            new TenantSaveChangesInterceptor(tenants));
}
\end{sourcecode}

With
\code{Add\allowbreak{}Db\allowbreak{}Context(\allowbreak{}(\allowbreak{}sp,\allowbreak{} o)
=> .\allowbreak{}.\allowbreak{}.\allowbreak{})}, \code{sp} is the scope of the current
request, because EF builds options per scope by default. The tenant provider and therefore
the interceptors see the right tenant. The registration lives in the core package because it
is the same for every provider (chapter~\ref{ch:what-belongs-in-the-core}). For the bridge
in the other direction, an interceptor can read services from the \emph{internal} provider
via
\code{event\allowbreak{}Data.\allowbreak{}Context.\allowbreak{}Get\allowbreak{}Service<\allowbreak{}T>(\allowbreak{})};
the connection interceptor gets the setting name from the singleton options that way
(chapter~\ref{ch:interceptors}).

\begin{important}[DbContext pooling]
  With \code{Add\allowbreak{}Db\allowbreak{}Context\allowbreak{}Pool}, contexts are reused
  and receive their options only once. Scoped interceptors from the first scope would be
  reused for all later requests, with the tenant of the first request. For pooling, the
  tenant provider must be a singleton that determines the tenant anew on every access, for
  example via \code{IHttp\allowbreak{}Context\allowbreak{}Accessor}.
  Chapter~\ref{ch:using-the-plugin-in-asp} shows this variant, and the same applies to
  \code{Add\allowbreak{}Db\allowbreak{}Context\allowbreak{}Factory}, whose factory is a
  singleton by default.
\end{important}
